\chapter{Model, Studien, Surrogate und Visualisierung}
\label{api:workflows}\index{Model}\index{Study}\index{Surrogat}
\begingroup\raggedright
\section{Model und Preflight}
\texttt{Model(mesh=..., document=...)} nimmt einen unabhängigen Dokumentstand,
\texttt{from\_file(path)} liest TOML und relative Mesh-/CSV-Pfade.
\texttt{to\_document()} liefert eine getrennte JSON-kompatible Kopie.
Regionen, Materialien, Einheiten, Rand-/Last-Zuordnungen, Lagerung/Gauge und
Elektroden werden vor Assembly geprüft; \texttt{solve()} erzeugt ModelResult
mit Physikname, passendem field\_result und RunProvenance. Schema 1 stationär,
Schema 2 transiente Wärme, Schema 3 Piezo-Moden/Newmark/Harmonik.
Keine implizite Schema-Migration, Remoteausführung oder automatische Vernetzung.
Schema-/Pfadfehler erhalten ModelSchemaError; physische Solverfehler bleiben
diagnostische Unterklassen. Interne PiezoDynamicModel-/TransientHeatModel-
Adapter sind keine zusätzlichen stabilen Anwender-Einstiegspunkte.
Minimalbeispiele: \nolinkurl{examples/thermopiezoelectric/case.toml},
\nolinkurl{examples/transient_heat/model.toml},
\nolinkurl{examples/piezo_dynamics/model.toml}; Tests lesen die Produktionsparser.

\section{Stationäre Parameterstudien und Mengen}
\texttt{Study(manifest, root=...)} beziehungsweise \texttt{from\_file}
hält detached Schema-1-Daten. Parameter: Name, vorhandener Zielpfad,
passende SI-Einheit und endliche eindeutige Werte. Ganze Skalar-/Vektor-/Tensor-
Eigenschaften, keine beliebigen Komponenten oder Ausdrücke. Parameter bilden
ein kartesisches Produkt in Deklarationsfolge, letzter Parameter am schnellsten.
Model/Netz müssen lokal innerhalb des Inputroots bleiben. \texttt{resolve()}
liefert detached Fälle ohne Solve; IDs hashen benannte Werte/Einheiten/Ziele,
Fingerprints zusätzlich effektives Modell, Meshbytes, Solver und Versionen.
Geänderte Inputs invalidieren Übernahme, auch wenn Fall-ID gleich bleibt.

\texttt{run(output, max\_cases=None)} arbeitet seriell mit Lock, persistiert
Versuche/Fehler und übernimmt nur hash-/identitätsgeprüfte Erfolge. Ein
Prozessabbruch hinterlässt Wiederaufnahmezustand; Teilversuche werden nicht
als erfolgreich wiederverwendet. \texttt{max\_cases} begrenzt neue Versuche,
nicht bereits erfolgreiche Übernahmen. Kein versteckter Multiprocessing-Dienst.
Minimalbeispiel: \nolinkurl{examples/parameter_study/run.py}; Heat-/Piezo-
Lebenszyklusfälle: \nolinkurl{examples/study_recipes/run.py}.

\texttt{Quantity} definiert eindeutigen Namen, metric/point/cell-Quelle, Feld,
SI-Einheit, fachliche Definition, Reduktion, Komponente und optionale Indizes.
Reduktionen sind ungeglättete ungegewichtete Rohwertreduktionen, keine
Flächen-/Volumenintegrale. \texttt{scalar} verlangt genau einen Wert;
vector-Komponente über Metadatenname oder norm. Optionale absolute Toleranz
bewertet den Betrag in derselben SI-Einheit, keine relativen Null-Bilanzen.
\texttt{export\_study\_table} schreibt neue CSV/JSON/Markdown-Verzeichnisse,
\texttt{read\_study\_table} prüft Schema, CSV-Hash und CSV-/JSON-Konsistenz.
Alle Fälle bleiben sichtbar, fehlende Werte samt Ursache; keine Solverausführung
beim Lesen. Export ersetzt kein Netzkonvergenz-/Messvalidierungsurteil.

\section{Inspectierbare polynomiale Surrogate}
\texttt{Variable} definiert Name/SI-Einheit/Definition.
\texttt{SurrogateData(X[n,p],Y[n,q],case\_ids,parameters,targets,provenance=...)}
kopiert endliche reelle Matrizen, eindeutige Fälle und Pflichtprovenienz auf
readonly Speicher. \texttt{from\_table} nutzt überprüfte Mengen und
prüft Fall-/Spalten-/Einheitenkonsistenz. Kein complex/bool/NaN oder stiller
Ausschluss fehlgeschlagener Fälle. DataSplit bindet getrennte Trainings-/Test-
IDs; \texttt{split\_cases} hält identische Eingabepunkte gruppiert und seeded.
Testdaten bestimmen weder Skalierung noch Modellfit.

\texttt{fit\_surrogate}: expliziter Grad/Ridge, Trainingsskalierung,
polynomiale Basis und prüfbare Koeffizienten. Unterbestimmter Rang,
inkonsistente IDs, Leakage und unzulässige Daten sind Fehler.
PolynomialSurrogate liefert Prediction mit Werte[n,q], Zielmetadaten und
Extrapolationsmaske[n]. Außerhalb der Trainingsbox wird standardmäßig
abgelehnt; erlaubte Extrapolation bleibt markiert. Trainings-/Test-MAE,
RMSE und Maximalfehler bleiben pro Ziel/SI-Einheit getrennt.
Speichern ist atomar, standardmäßig ohne Überschreiben; Laden prüft Schema,
Hash, Formen, Training/Test-Provenienz und tatsächliche Artefaktstruktur.
Approximationsfehler ist keine experimentelle Eingabeunsicherheit oder
FEM-Netzfehlergrenze. Beispiel: \nolinkurl{examples/surrogate_heat/run.py}.

\section{Optionale Bilder und Wärmeindikatoren}
\texttt{plot\_field(mesh, values, ...)} verarbeitet reelle rohe Tet4-
Feldarrays mit \texttt{association=node/cell}, explizitem Namen/SI-Einheit,
aktiven Knoten und optionaler xyz-/Voigt-Komponente oder vector magnitude.
Tet10 und komplexe Felder werden abgelehnt. FieldPlot enthält owned Figure,
Axes, Artist, Außenflächen-Indizes und tatsächliche Farbwerte.
Knotenfarben sind arithmetische Flächenmittel nach Komponenten-/Normauswertung;
Zellfarben kommen ungeglättet aus der owning volume cell. Keine automatische
von-Mises-Größe. Explizite endliche wachsende Farblimits oder dokumentierte
Erweiterung eines konstanten automatischen Bereichs.
\texttt{save\_jpeg(plot, path, dpi=150, overwrite=False)} schreibt JPEG in
ein bestehendes Verzeichnis; Pixelgröße ist figsize in Zoll mal DPI.
Matplotlib ist ein separates Extra, erst bei Plotaufruf erforderlich.
Aufrufer schließt Figuren; mutable Plotobjekte sind keine frozen Resultate.
Kein pyplot-Registrieren, show, globaler Style- oder Backendwechsel.
Beispiel: \nolinkurl{examples/field_visualization/run.py}.

\texttt{heat\_indicator}/\texttt{mark\_bulk} sind Forschungs-API für stationäre
Tet4-Wärme: Residualbeiträge, Flächensprünge und deterministische Markierung.
Kein garantierter globaler Fehler, kein Remeshing und kein Tet10-Transfer.
Definition/Skalierung stehen in Doc/43; niedrigste Assembly-/Element-/Physik-
Helfer bleiben genauso Forschung oder interne Erweiterungsgrenzen.
\endgroup
\begingroup
\raggedright
\section{Unsicherheits-API und Artefakte (M9 / 64)}
\index{SamplingPlan}\index{UncertaintyStudy}\index{Sensitivität}
Modul \texttt{fem\_toolbox.uncertainty}; der Paketroot bleibt unverändert.

\paragraph{Normal} \texttt{Normal(mean, standard\_deviation)}. Endliche reelle SI-Skalare,
positive Standardabweichung, unbeschränktes Normalgesetz, unveränderlicher Datensatz.
\paragraph{Uniform} \texttt{Uniform(lower, upper)}. Endliches wachsendes SI-Intervall mit endlicher
Breite, unveränderlicher Datensatz. Keine automatische Rechtfertigung des Gesetzes.
\paragraph{UncertainParameter} Name (ASCII), vorhandener skalarer M8-Zielpfad,
SI-Einheit, Normal-/Uniformgesetz und Pflichtbegründung. Frozen-Datensatz.
\subsubsection*{SamplingPlan}
\begin{lstlisting}
SamplingPlan(parameters, *, size, seed, correlation=None)
\end{lstlisting}
$n\ge2$, nichtnegativer Integer-Seed, eindeutige Namen/Ziele. Reelle
$R\in\mathbb R^{p\times p}$ symmetrisch, PSD, Einserdiagonale; absolute
Rundungstoleranz $10^{-12}$. Nichtidentität nur für alle-normal; singulär erlaubt.
\texttt{sample()} liefert bytebasiertes unveränderliches float64$[n,p]$.
\texttt{to\_document()} liefert getrennte veränderbare JSON-Metadaten.
\subsubsection*{Stichprobenstatistik}
\begin{lstlisting}
summarize_samples(values, *, unit, numerical_error_bound,
                  quantile_probabilities=(.025,.5,.975))
\end{lstlisting}
Reeller endlicher Vektor $[n\ge2]$, explizite SI-Einheit. Budget nichtnegativ
oder \texttt{None}; geordnete Quantile in $[0,1]$, Default $(.025,.5,.975)$.
Mittel, unverzerrte Varianz (Einheit$^2$), linear interpolierte Quantile und
Präfixmomente; keine Konfidenz-/Fehlergrenze. Budget bleibt getrennt.
\subsubsection*{Lokale Sensitivität}
\begin{lstlisting}
local_sensitivity(evaluate, *, center, steps)
\end{lstlisting}
Reelle endliche Vektoren $[p]$, positive unterscheidbare Perturbationen;
genau $2p$ skalare Callback-Ausgaben. Lokaler zentraler Gradient;
Einheit Ausgabe/Eingabe. Callbackfehler propagieren.
\subsubsection*{UncertaintyStudy}
\begin{lstlisting}
UncertaintyStudy(plan, *, model, root)
run(output, *, quantities, numerical_error_bounds,
    sensitivity_center=None, sensitivity_steps=None)
\end{lstlisting} Relatives stationäres TOML-Modell
mit M8-Pfad-/SI-Vertrag. Die Methode \texttt{run}
fordert eindeutige Quantity-Definitionen und Budget/None für jedes Ziel;
Sensitivitätszentrum/-schritte müssen gemeinsam $[p]$ sein.

Bool/complex/NaN, falsche Formen, doppelte Identitäten und unzulässige
Verteilungen erzeugen \texttt{ValueError}. M8-Schemafehler propagieren;
physisch ungültige Fälle bleiben als Ergebnisstatus sichtbar. I/O-/Hashfehler
propagieren, konkurrierende Writer erzeugen \texttt{FileExistsError}.

Schema 1: \texttt{sampling.json} hält Gesetze, Begründungen, Reihenfolge, $R$,
PCG64/Transformationsmethode, $n$, Seed, NumPy-/SciPy-Versionen fest.
\texttt{uncertainty.json} enthält Samplingprovenienz, Little-Endian-float64-
Stichprobenhash, Quantity-Verträge, gemeinsame Zeilen mit M8-Provenienz,
Katalog-/Tabellenpfaden und Hashes, Statistik sowie Sensitivitätsbelege.
\texttt{complete=false} unterdrückt Statistik, sobald ein Fall oder eine Größe
fehlt; keine still ausgewählte Restmenge. Misslungene zusätzliche Perturbationen
setzen \texttt{sensitivity\_status=incomplete}. Geänderte Modelle invalidieren
M8-Ergebnisübernahme; geänderte Samplingpläne brauchen neue Ausgaben.
\texttt{.uncertainty.lock} schützt Writer. Nach hartem Abbruch vor manueller
Lock-Entfernung einen verbleibenden Prozess ausschließen.

Die neue öffentliche Modulgruppe und der Quellenvertrag sind in der
Abdeckungsmatrix enthalten. Signaturdrift dieses Moduls wird durch den AST-Vertrag
\nolinkurl{Doc/latex/uncertainty-api.json} geprüft (Funktionen, Felder und Methoden).
Das vollständige API-Inventar aller Module folgt in \#209. Das Benutzerrezept und die unabhängigen
analytischen Tests gehören bereits zum Pflichtumfang von \#64.

\endgroup

